% ECONOMICS_PROBLEM: {"id": "C73-10", "title": "Ordinary uniform approximate equilibria in quitting games", "jel_code": "C73", "source_id": "OP-0591"}
% FIRST_POSTED: 2026-10-09
% ATTEMPT_STATUS: unresolved
% ENTRY_KIND: research_attempt
% !TeX program = pdflatex
\documentclass[11pt,a4paper]{article}
\usepackage[T1]{fontenc}
\usepackage[utf8]{inputenc}
\usepackage{lmodern}
\usepackage[margin=24mm,headheight=15pt]{geometry}
\usepackage{amsmath,amssymb,amsthm,mathtools}
\usepackage{booktabs,longtable,array,enumitem}
\usepackage{microtype}
\usepackage{xurl}
\usepackage[hidelinks]{hyperref}
\usepackage{fancyhdr}
\newcommand{\E}{\mathbb E}
\newcommand{\R}{\mathbb R}
\newcommand{\N}{\mathbb N}
\newcommand{\ind}{\mathbf 1}
\DeclareMathOperator{\Reg}{Reg}
\DeclareMathOperator{\supp}{supp}
\newtheorem{theorem}{Theorem}[section]
\newtheorem{lemma}[theorem]{Lemma}
\newtheorem{proposition}[theorem]{Proposition}
\newtheorem{corollary}[theorem]{Corollary}
\theoremstyle{definition}
\newtheorem{definition}[theorem]{Definition}
\theoremstyle{remark}
\newtheorem{remark}[theorem]{Remark}
\setlength{\parindent}{0pt}
\setlength{\parskip}{4pt plus 1pt}
\setlength{\emergencystretch}{2em}
\setcounter{tocdepth}{1}
\allowdisplaybreaks[2]
\pagestyle{fancy}
\fancyhf{}
\fancyhead[R]{\small Open Problems Econ}
\fancyfoot[C]{\thepage}
\renewcommand{\headrulewidth}{0.3pt}
\hypersetup{pdfauthor={GPT 6 Astra Pro},pdfsubject={Checked partial results and explicit remaining gaps}}

\fancyhead[L]{\small\texttt{C73-10} -- Research attempt}
\title{Research attempt: \texttt{C73-10}\\[5pt]\large Ordinary approximate equilibria in quitting games}
\author{GPT 6 Astra Pro}
\date{9 October 2026}
\begin{document}
\maketitle
\noindent\textbf{Scope of this report.} The main target remains UNRESOLVED. Fully proved subclaims and counterexamples are identified locally. Computational checks reported here are supplementary to the proofs. Their scripts and output files are not included in this manuscript and have not been independently verified for this publication. No novelty or priority claim is made.\par
\section{FORMAL RESTATEMENT}
Let $n\geq2$ and $r:2^{[n]}\to\mathbb R^n$ be any finite payoff table, with $c=r(\varnothing)$. At each date $t\in\mathbb N_0=\{0,1,\ldots\}$, every still-active player independently chooses continue or quit. The first nonempty quitting set $J$ causes absorption and payoff $r(J)$ at that date and forever after; before it, stage payoff is $c$. If nobody ever quits, terminal payoff is $c$. No public correlation device is allowed.

Equivalently, player $i$ chooses a probability law $\mu_i$ on $K=\mathbb N_0\cup\{\infty\}$, independently across players. The first finite minimum of the times determines $J$. Denote the resulting terminal expected payoff by $\Gamma_i(\mu)$. Define
\[
 \Reg_r(\mu)=\max_i\sup_{\nu_i}
 [\Gamma_i(\nu_i,\mu_{-i})-\Gamma_i(\mu)],\qquad
 g(r)=\inf_\mu\Reg_r(\mu).
\]
All payoffs are bounded automatically because the table is finite. Let
\[
 D=\max_i\left(\max_Jr_i(J)-\min_Jr_i(J)\right).
\]
The terminal assertion is $\forall n\geq2\ \forall r\ \forall\varepsilon>0\ \exists\mu$ with regret at most $\varepsilon$, equivalently $g(r)=0$ for every table.

The uniform assertion requires one profile $\mu^\varepsilon$, a finite $H_\varepsilon$, and $\lambda_\varepsilon>0$ such that every player and every unilateral behavioral deviation has gain at most $\varepsilon$ both for all $H\geq H_\varepsilon$ in the expected $H$-stage average and all $0<\lambda\leq\lambda_\varepsilon$ in $\sum_{t\geq0}\lambda(1-\lambda)^t$-weighted stage rewards.

For $k\geq1$, let $\mathcal G_{n,k}$ consist of product laws whose marginals are empirical laws of $k$ points of $\{0,\ldots,nk-1,\infty\}$, with repetitions allowed, and put
\[
 m_k(r)=\min_{\mu\in\mathcal G_{n,k}}\Reg_r(\mu).
\]
The problem statement also asks whether $m_k(r)\leq(2n-1)D/k$ universally, and separately asks about stationary, finite-memory and subgame-perfect refinements. Regret in this definition always tests \emph{unrestricted} deviations, not just deviations in the grid.

\paragraph{Hygiene and source availability.}
We independently prove the grid reduction and the equivalence of the terminal and uniform existence assertions. If $D=0$, every profile is an exact equilibrium and all asserted regret bounds are zero. Simultaneous quitting and the no-quitting payoff remain part of every construction.

\section{QUICK LITERATURE/CONTEXT CHECK}
The foundational quitting-game paper distinguishes ordinary strategic issues from particular stationary constructions \cite{c7310-sv}. Absorption-path work accommodates both discrete jumps and continuous hazards \cite{c7310-ap}. A 2026 APS paper characterizes a subset of subgame-perfect approximate payoffs through absorption paths with exactly one active player; that restriction is not a proof of the general ordinary-equilibrium assertion \cite{c7310-aps}. The present results neither introduce a sunspot device nor identify stationary nonexistence with ordinary nonexistence. Sources checked through 9 October 2026.

\section{ATTACK PLAN}
\paragraph{Proof strategies.}
Convert behavioral strategies into independent stopping-time laws; quantize each marginal while controlling every possible deviation; then transfer terminal regret to long finite and discounted evaluations. This gives a quantitative finite-certificate reformulation but leaves a universal grid inequality to prove.

\paragraph{Disproof strategies.}
Search exact small tables for positive grid regret; seek a table whose regret exceeds the rounding allowance; and analyze a mismatch construction that destroys exact equilibrium. The mismatch construction succeeds only against the \emph{stronger exact-existence} claim.

\paragraph{Landscape and tools.}
This is discontinuous infinite-action game theory. Hazard-rate representation removes irrelevant histories; midpoint quantiles control distribution functions; bounded variation of the winner/tie payoff controls rounding error; hybrid replacement preserves independent randomization; order-preserving time compression makes the search finite; affine payoff dependence reduces best responses to pure times; and absorption-lag bounds transfer terminal incentives uniformly over horizons.

\section{WORK}
\subsection{Stopping times and a quantitative grid reduction}
\begin{lemma}[Behavioral equivalence and pure best responses]
Every behavioral strategy profile induces the same outcome distribution as independent laws on $K$, and conversely. Against fixed opponents, the best-response payoff is the supremum over deterministic times in $K$.
\end{lemma}
\begin{proof}
At a survived date the only previous joint actions are all continue. If player $i$ quits with conditional probability $p_i(t)$, define
\[
 \mu_i(t)=p_i(t)\prod_{s<t}(1-p_i(s)),\qquad
 \mu_i(\infty)=\prod_{s\geq0}(1-p_i(s)).
\]
Independent randomization produces exactly the independent planned times with these marginals. Conversely, set the hazard to the mass at $t$ divided by the remaining survival mass, choosing it arbitrarily if the denominator is zero. This reconstructs every marginal. A unilateral behavioral deviation has the same representation; no history with a past quit can affect the outcome. Finally, the payoff from a mixed own time is the integral of the deterministic-time payoff, hence cannot exceed its supremum; point masses attain every deterministic-time value. Attainment of the supremum itself is not assumed.
\end{proof}

\begin{lemma}[Midpoint quantiles]
Every law $\mu$ on $K$ has a $k$-point empirical law $\nu$ such that, at each finite $t$,
\[
 |\mu(\{0,\ldots,t\})-\nu(\{0,\ldots,t\})|\leq\frac1{2k}.
\]
Replacing one player's law by this empirical law changes any player's expected payoff by at most $D/k$, uniformly in all the other laws.
\end{lemma}
\begin{proof}
Order $\infty$ after all integers and choose quantiles at $(2\ell-1)/(2k)$ for $\ell=1,\ldots,k$, sending a quantile to $\infty$ when the finite cumulative probabilities never reach its level. At a finite cumulative probability $F(t)$, the number of chosen quantiles at or below $t$ differs from $kF(t)$ by at most $1/2$. This proves the bound, including atoms and ties. The left limit $F(t-1)$ has the same bound, with $F(-1)=0$.

Fix deterministic times of all players except the law being replaced. If their earliest time is finite, say $t$, the payoff to any fixed player, as a function of the replaced time, has three values $a,b,c'$: before $t$, at $t$, and after $t$. Its expectation changes by
\[
 (a-b)\bigl(F_\mu(t-1)-F_\nu(t-1)\bigr)
 +(b-c')\bigl(F_\mu(t)-F_\nu(t)\bigr).
\]
Each coefficient has absolute value at most $D$. The absolute change is therefore at most $D/k$. If all other times are infinite, there are only the finite-time and infinite-time values; the change is at most $D/(2k)$, using the limit of the finite cumulative probabilities. Integrate the deterministic-other-times bound against their independent laws to obtain the uniform claim.
\end{proof}

\begin{theorem}[Grid approximation with unrestricted-deviation certification]
For every payoff table and every $k\geq1$,
\[
 \boxed{\quad g(r)\leq m_k(r)\leq g(r)+\frac{(2n-1)D}{k}.\quad}
\]
Consequently $m_k(r)\to g(r)$, and
\[
 g(r)=0\ \Longleftrightarrow\
 \forall k\geq1:\quad m_k(r)\leq\frac{(2n-1)D}{k}.
\]
\end{theorem}
\begin{proof}
Round all marginals of an arbitrary profile $\mu$ by the lemma, obtaining $\nu$. Replacing them one at a time changes on-profile payoffs by at most $nD/k$. Keeping an arbitrary deviator fixed while replacing the other $n-1$ marginals changes its deviating payoff by at most $(n-1)D/k$. Thus
\[
 \Reg_r(\nu)\leq\Reg_r(\mu)+(2n-1)D/k.
\]
The estimate holds simultaneously for all deviations, so taking their supremum is justified.

Across the $n$ rounded marginals there are at most $nk$ distinct finite support times, say $z_1<\cdots<z_M$. Compress them to $0,\ldots,M-1$, preserving $\infty$. On-profile outcomes, including all ties, are unchanged. Every deterministic deviation at a compressed time $t<M$ is matched by the original time $z_{t+1}$; every finite deviation $t\geq M$ is matched by $z_M+1$; and an infinite deviation is matched by $\infty$. If $M=0$, use any finite original time for a finite deviation. These correspondences preserve the first quitting set against the opponents. Compression can remove opportunities to quit in an old gap, but it adds no opportunity with a new payoff. By the pure best-response lemma, compressed regret cannot increase. The resulting profile lies in $\mathcal G_{n,k}$.

Choose $\mu$ with regret at most $g(r)+\delta$, possible by the definition of an infimum, and let $\delta\downarrow0$ in the resulting bound on the finite minimum $m_k$. The reverse inequality $g\leq m_k$ holds because the grid is a subset of all profiles. The convergence and equivalence follow. No existence of an infimum-attaining profile was assumed.
\end{proof}

\paragraph{An exact finite certificate and a necessary search range.}
For a grid profile, it is enough to check deterministic deviation times
\[
 0,1,\ldots,nk,\infty.
\]
The extra time $nk$ is essential: it quits strictly after all finite opponent support, and is not in the profile-support grid. Every larger finite time has the same terminal effect. Each marginal has $\binom{(n+1)k}{k}$ possible count vectors, so exhaustive rational search is finite. Because regret is always at most $D$, a strict violation of the universal grid inequality is impossible when $k\leq2n-1$. Any counterexample certificate must have $D>0$ and $k\geq2n$.

\subsection{Terminal and uniform existence really are equivalent}
\begin{theorem}[Common-profile transfer]
For each fixed table, $g(r)=0$ if and only if the joint finite-horizon and discounted uniform approximate-existence assertion in Section 1 holds.
\end{theorem}
\begin{proof}
Suppose first $g(r)=0$ and $D>0$. Given $\varepsilon>0$, choose $k$ so that $(2n-1)D/k\leq\varepsilon/2$. The grid theorem supplies a profile $\mu$ with terminal regret $\delta\leq\varepsilon/2$ and all finite support before some integer $T\geq1$.

On this profile, any absorption occurs before $T$, so replacing the finite or discounted payoff by the terminal payoff makes an error at most $DT/H$ or $D\lambda T$, respectively. This follows pathwise: at most $T$ initial rewards differ for the average, and their discounted mass is at most $1-(1-\lambda)^T\leq\lambda T$.

Fix a deterministic deviating time $t$. For $t\leq T$, the same bounds apply against the grid opponents. For $t>T$, separate the event that some opponent has a finite time from the event that all opponents have time $\infty$. On the first event, the payoff absorbs before $T$ and the same error bound applies. On the second, the deviator alone switches from $c_i$ to $r_i(\{i\})$ at time $t$. Its finite-horizon payoff is a convex combination, with coefficient $(1-t/H)_+$, of these two values; its discounted payoff has coefficient $(1-\lambda)^t$. The corresponding full expected payoff is therefore, up to the early-absorption error, a convex combination of the two terminal pure-deviation payoffs obtained by choosing time $T$ and time $\infty$. Both are bounded above by the terminal best-response supremum.

The argument is uniform in the deterministic time, hence, by integrating, in every deviating law. Subtracting the on-profile payoff bound gives
\[
 \Reg^H(\mu)\leq\delta+2DT/H,
 \qquad \Reg^\lambda(\mu)\leq\delta+2D\lambda T.
\]
Choose $H\geq4DT/\varepsilon$ and $\lambda\leq\min\{1/2,\varepsilon/(4DT)\}$. The same profile then works for all these horizons, all these discount rates, and all deviations. The case $D=0$ is immediate.

Conversely, take a profile supplied by uniform existence at accuracy $\varepsilon$. For this fixed profile and any fixed unilateral deviation, the stage averages converge pathwise to their respective terminal payoffs: they are eventually constant after absorption, or identically $c$ if no one quits. Bounded convergence passes each expected average to the terminal expectation. Taking the large-horizon limit of the incentive inequality gives terminal gain at most $\varepsilon$ for this arbitrary deviation. Since this holds for every deviation, terminal regret is at most $\varepsilon$. No interchange of a limit with a supremum is used. Let $\varepsilon\downarrow0$ in the existence assertion to obtain $g(r)=0$.
\end{proof}

\subsection{An exact nonexistence example that does not refute approximation}
Consider two players with table
\[
\begin{array}{c|rrrr}
 J&\varnothing&\{1\}&\{2\}&\{1,2\}\\\hline
 r_1(J)&0&1&1&0\\
 r_2(J)&1&0&0&1
\end{array}
\]
Terminal payoffs are $(\ind_{\{T_1\ne T_2\}},\ind_{\{T_1=T_2\}})$, including the case $T_1=T_2=\infty$.

\begin{theorem}[Exact grid minimum for the mismatch game]
This game has no exact Nash equilibrium, but $g(r)=0$. For every $k\geq1$,
\[
 \boxed{\quad m_k(r)=\frac{\lceil k/2\rceil}{k^2}.\quad}
\]
In particular the general $O(1/k)$ grid-approximation order cannot be improved to $o(1/k)$ uniformly over all tables with $g(r)=0$.
\end{theorem}
\begin{proof}
For arbitrary laws let $a=\max_t\mu_1(t)$ and $c_0=\sum_t\mu_1(t)\mu_2(t)$. A probability on a countable set has a strictly positive largest atom: its supremum is positive, and if it were not attained infinitely many atoms would exceed half that positive supremum, contradicting summability. Thus $a>0$. Also $\inf_t\mu_2(t)=0$ since $K$ is infinite. Player 1's best-response supremum is 1, while player 2's best-response maximum is $a$. Their regrets are consequently $c_0$ and $a-c_0$, giving
\[
 \Reg_r(\mu)=\max\{c_0,a-c_0\}\geq a/2>0.
\]
This rules out exact equilibrium but not approximate equilibrium.

On the denominator-$k$ grid, if $a\geq2/k$, regret is at least $1/k$. Otherwise $\mu_1$ is uniform on exactly $k$ distinct points, $a=1/k$, and $c_0=\ell/k^2$ for an integer $0\leq\ell\leq k$. Hence regret is at least
$\min_{\ell}\max\{\ell,k-\ell\}/k^2=\lceil k/2\rceil/k^2$.
This bound is at most $1/k$, so it also covers the first case.

To attain it, let $\mu_1$ be uniform on $0,\ldots,k-1$. Let $\mu_2$ put $\lfloor k/2\rfloor$ of its $k$ empirical points at time 0 and all remaining points at time $k$. These times lie in $\{0,\ldots,2k-1\}$, including when $k=1$. Then $c_0=\lfloor k/2\rfloor/k^2$ and the claimed regret follows. Letting $k\to\infty$ proves $g(r)=0$. For even $k$ the exact grid minimum is $1/(2k)$, proving the asserted rate obstruction.
\end{proof}

\section{VERIFICATION}
The reported computation exhaustively evaluates all 256 binary-payoff two-player quitting tables on the $k=4$ grid, with unrestricted best responses reduced to $0,\ldots,8,\infty$. There are 495 marginal count vectors and 245,025 product profiles per table. Exact integer numerators show that 254 tables have $m_4=0$ and the mismatch table and its player-payoff reversal have $m_4=1/8$. None violates the allowance $3D/4$. The computation also verifies 810 finite midpoint-quantile distribution-function inequalities. Finite search is not a search over all real tables or all $k$.

Pseudocode is: enumerate payoff tables; enumerate both empirical count vectors; calculate current payoffs and each pure-time deviation by integer products of counts; minimize the maximum gain; compare the minimum with the rational rounding threshold. A grid search omitting time $nk$ would not be an adequate certificate. The analytic mismatch proof checks every real distribution, not only enumerated ones.

Compression preserved the order and ties of finite support times and retained the infinite atom. The common-profile transfer controlled late deviations by a convex combination rather than falsely asserting uniform convergence over all stopping times. No result here implies subgame perfection or an accuracy-independent finite-memory bound. The mismatch profile constructed for a fixed $k$ can be implemented with a finite date counter, but this says nothing about general games.

\section{FINAL}
\textbf{LABEL: UNRESOLVED} for universal ordinary approximate equilibrium existence and its uniform equivalent.

\paragraph{Strongest checked partial results.}
The exact inequalities $g\leq m_k\leq g+(2n-1)D/k$ and the terminal--uniform equivalence are fully proved. The mismatch game has no exact equilibrium but has $m_k=\lceil k/2\rceil/k^2$ and therefore arbitrarily accurate ordinary and uniform equilibria. These are solved reductions and a solved example, not a universal existence proof.

\paragraph{Exact first gap.}
For an arbitrary table $r$, prove $m_k(r)\leq(2n-1)D/k$ for every $k$, or exhibit one table and one $k$ with a rigorously certified strict reverse inequality. The approximation theorem bounds $m_k$ in terms of the unknown $g$ and does not supply this missing universal bound.

\paragraph{Three next moves.}
(1) Target four-player signed tables on grids with $k\geq8$, where the allowance is no longer automatically larger than the whole payoff range. (2) Develop exact lower-bound certificates for all empirical product profiles, rather than checking a heuristic stationary family. (3) Seek an absorption-path construction that permits both ties at discrete jumps and simultaneous continuous hazards and proves the missing grid bound.

\paragraph{Minimal counterexample structure.}
A valid approximate-existence counterexample would have positive $g(r)$, not just a missing exact or stationary equilibrium. It must retain its full simultaneous-quitting and no-quitting table. Any finite-grid certificate must satisfy $D>0$ and $k\geq2n$ and test all unilateral times described above.

\begin{thebibliography}{9}
\bibitem{c7310-sv} E. Solan and N. Vieille. \emph{Quitting Games}. Mathematics of Operations Research 26 (2001), 265--285. \url{https://doi.org/10.1287/moor.26.2.265.10549}.
\bibitem{c7310-ap} G. Ashkenazi-Golan, I. Krasikov, C. Rainer, and E. Solan. \emph{Absorption paths and equilibria in quitting games}. Mathematical Programming 203 (2024), 735--762. \url{https://doi.org/10.1007/s10107-022-01807-6}.
\bibitem{c7310-aps} G. Ashkenazi-Golan, I. Krasikov, C. Rainer, and E. Solan. \emph{The APS approach for undiscounted quitting games}. International Journal of Game Theory 55 (2026), article 19. \url{https://doi.org/10.1007/s00182-026-00982-6}.
\end{thebibliography}

\end{document}
